\SourceLeaf{129}
und Leichtigkeit entnommen werden könne, schon vorhanden sind. Jedenfalls aber würden doch die ohne Zweifel aufbewahrten 105 Jahresrechnungen dazu dienen können. Erst nach eigner näherer Einsicht in die vorhandenen Papiere würde ich aber mich erklären können, ob und in welchem Maasse ich meine eigne Beihülfe zu dieser Extraction zusagen kann.

Für jeden einzelnen in diesem Buche aufzuführenden Theilnehmer lässt sich aus den rubricirten Daten berechnen, wie viel er der Gesellschaft und wie viel seinen Relicten diese bei den gegenwärtigen Sätzen geleistet haben; und was bei bestimmtem Zinsfusse der Geldwerth davon auf die Zeit seines Eintritts reducirt gewesen sein würde. Natürlich ist dies bei den einzelnen sehr verschieden, bei einigen positiv, bei andern negativ; aber nach den ewigen Gesetzen ist der \textit{Mittelwerth} aus einer grossen Menge ein Element das als Mittelwerth wieder für die Zukunft zum Grunde gelegt werden kann, wo keine Ursache ist, wesentliche Aenderungen der allgemeinen Verhältnisse vorauszusetzen. Diese Tabelle selbst wird hierüber schon eine lehrreiche Indication geben können, wenn man das Ganze gruppirt, und z. B. den Mittelwerth der ersten Hälfte mit dem für die zweite, oder das erste, zweite und dritte Drittel mit einander vergleicht. Es wird sich so herausstellen, wie gross der Geldwerth der Verbindlichkeit der Gesellschaft ist, der ihr durch den Eintritt eines neuen Mitgliedes \textit{durchschnittlich} zuwächst, und wenn man, nach dem was schon oben bemerkt ist, zugleich eine plausible Annahme für die Durchschnittszahl der jährlich zutretenden genommen hat, so lässt sich, für bestimmten Zinsfuss, berechnen, wie gross das Capital sein muss, dessen Zinserträgniss, diese auf immer fortlaufende Verbindlichkeit III. decken kann. Alle drei Capitale für I. II. und III. zusammengerechnet und mit dem wirklichen Vermögen der Gesellschaft verglichen, werden dann so genau wie es nach der Natur des Gegenstandes möglich ist zeigen, ob bei der jetzigen Einrichtung ihre Stabilität \textit{mehr} als gesichert ist, oder nur eben zureichend, oder aber ob ihre Instabilität daraus hervorgeht, und also mit Entschiedenheit früh oder spät ihr Untergang erwartet werden müsse, und demgemäss würden dann die geeigneten Maassregeln in den verschiedenen Fällen zu erwägen und einzuleiten sein.

Dies sind die Hauptzüge des Planes, nach welchem meiner Meinung nach eine gründliche Prüfung und Aufstellung einer Bilanz ausgeführt werden könnte und müsste. Es ist eine bedeutende Arbeit, der ich mich aber, wenn es gewünscht wird, gern unterziehen werde. Dass diese Skizze nur ungefeilt und flüchtig niedergeschrieben hier vorgelegt ist, wird man mit der Kürze der Zeit entschuldigen. Wird eine solche Arbeit jetzt ausgeführt, so bleibt es jedenfalls wie schon oben bemerkt ist, dringend wünschenswerth, dass in Zukunft nach gewissen Zeitfristen immer wieder eine neue Bilanz gezogen werde, und dies würde dann viel leichter als das erstemal werden, wenn ein solches Buch wie ich oben erwähnt habe mit allen \textit{Zeitpunktsrubriken} wenigstens von jetzt an angelegt und regelmässig vervollständigt und fortgesetzt würde.

Ich will nun auch noch mein Votum über ein paar andere Punkte, die in den andern Abstimmungen berührt sind, beifügen.

Ich bin nicht dafür, dass die Bestimmung der Statuten, welche die nicht besoldeten Professoren ausschliesst, aufgehoben werde. Die Universität als Corporation müsste in Beziehung auf die Witwenkasse immer dringend wünschen, dass solche Fälle, wo einem bei Schule oder Kirche Angestellten der Professortitel beigelegt wird, sehr selten blieben. Von einem solchen Fall aber abgesehen, wird einer, der gar keine Besoldung und kein Vermögen hat, nicht leicht so unbesonnen sein, eine Verheirathung einzugehen, und also überhaupt die ganze Bestimmung selten vielleicht nie von Wirkung sein. Möglicherweise könnte aber ein unbesoldeter Professor, der sich selbst dem Tode nahe fühlend eine Braut hätte, welche er sonst vor Erlangung einer Besoldung gewiss noch nicht geheirathet hätte, falls ihm der Eintritt in die Witwenkasse offen stände, dadurch versucht werden, durch eine schnelle Copulation der Witwenkasse eine Last aufzubürden.

So lange \textit{bona fide} gehandelt wird, müssen vielmehr die unbesoldeten Professoren jene Clausel als
